% Cover
\thispagestyle{empty}
\begin{tikzpicture}[remember picture,overlay]
  \fill[Navy] (current page.north west) rectangle ([yshift=-62mm]current page.north east);
  \fill[Blue] ([yshift=-62mm]current page.north west) rectangle ([yshift=-66mm]current page.north east);
  \node[anchor=north,text=white,align=center] at ([yshift=-10mm]current page.north)
    {\Large 電気通信大学大学院};
  \node[anchor=north,text=white,align=center] at ([yshift=-20mm]current page.north)
    {\large 基盤理工学専攻 一般入試対策};
  \node[anchor=north,text=white,align=center] at ([yshift=-33mm]current page.north)
    {\fontsize{28}{34}\selectfont\bfseries 基礎数学 予想模試};
  \node[anchor=north,text=white,align=center] at ([yshift=-51mm]current page.north)
    {\Large 2026年度実施試験 対応};
\end{tikzpicture}

\vspace*{62mm}
\begin{center}
\begin{minipage}{0.82\textwidth}
\begin{examinfo}
\centering
\begin{tabular}{>{\bfseries}rl}
推奨時間 & 45分 \\
満点 & 100点 \\
構成 & 線形代数・多変数解析・重積分・微分方程式 \\
難易度 & 本試験標準 - やや難 \\
\end{tabular}
\end{examinfo}
\end{minipage}
\end{center}

\vfill
\begin{center}
\begin{minipage}{0.82\textwidth}
\small
本冊子は、2019--2025年度の出題傾向を踏まえて作成したオリジナル模試である。近年の頻出分野である固有値・固有空間を軸に、出題間隔を考慮して変数変換による重積分と常微分方程式を組み合わせた。
\end{minipage}
\end{center}
\vspace{12mm}
\begin{flushright}
氏名：\rule{55mm}{0.4pt}\qquad 得点：\rule{20mm}{0.4pt} / 100
\end{flushright}
\newpage

% Instructions
\section*{受験上の注意}
\begin{examinfo}
\begin{enumerate}
  \item 制限時間は45分を目安とする。
  \item 電卓・参考書は使用しない。
  \item 結果だけでなく、主要な計算過程を記すこと。
  \item 行列の固有空間では、基底を列ベクトルで明示すること。
  \item 重積分では、変数変換・領域・ヤコビアンを必ず示すこと。
  \item 微分方程式では、一般解を求めた後に初期条件を適用すること。
\end{enumerate}
\end{examinfo}

\vspace{5mm}
\section*{配点と時間配分の目安}
\begin{center}
\begin{tabularx}{0.92\textwidth}{c X c c}
\toprule
大問 & 主題 & 配点 & 目安時間 \\
\midrule
1 & 固有値・固有空間・対角化・行列のべき & 40点 & 18分 \\
2 & 極値判定・変数変換による重積分 & 35点 & 16分 \\
3 & 1階・2階常微分方程式 & 25点 & 11分 \\
\bottomrule
\end{tabularx}
\end{center}

\vspace{6mm}
\begin{pointbox}
\textbf{得点目標}
\quad 80点以上：非常に良い \quad 65--79点：合格圏 \quad 50--64点：要復習 \quad 49点以下：基礎事項を再確認
\end{pointbox}

\vfill
\begin{center}
\Large\bfseries 問題は次ページから
\end{center}
\newpage
